\subsubsection{Circuit notation and gate composition}\label{ex:circuit-notation-and-gate-composition}

In this section, we will not introduce any complex exercises. Instead, we will explore how to approach circuit notation and compose gates in a circuit.

\highspace
\begin{flushleft}
    \textcolor{Green3}{\faIcon{long-arrow-alt-right} \textbf{Consecutive Gates}}
\end{flushleft}
Consider this circuit:
\begin{center}
    \begin{quantikz}
        \lstick{\ket{\psi}} & \gate{U_1} & \gate{U_2} & 
    \end{quantikz}
\end{center}

\noindent
Physically, this circuit applies the gate $U_1$ to the state $\ket{\psi}$, and then applies the gate $U_2$ to the resulting state. But mathematically we need \textbf{one equivalent matrix} to represent the entire circuit. This is done by multiplying the matrices of the gates in the order they are applied:
\begin{equation*}
    U_3 = U_2 U_1
\end{equation*}
Notice the order of multiplication: the \textbf{rightmost matrix acts first}.

\highspace
\textcolor{Green3}{\faIcon{question-circle} \textbf{Why?}} This rule comes from linear algebra. If $\ket{\psi}$ is a column vector, first we apply $U_1$ obtaining $U_1 \ket{\psi}$. Then we apply $U_2$ giving $U_2 \left(U_1 \ket{\psi}\right)$. Associativity of matrix multiplication allows us to write this as $\left(U_2 U_1\right) \ket{\psi}$, which is equivalent to applying the single gate $U_3 = U_2 U_1$ to the state $\ket{\psi}$.

\highspace
\begin{examplebox}[: Circuit notation and gate composition]
    Consider the following circuit:
    \begin{center}
        \begin{quantikz}
            \lstick{\ket{\psi}} & \gate{H} & \gate{Z} & 
        \end{quantikz}
    \end{center}
    Physically:
    \begin{enumerate}
        \item Hadamard
        \item Pauli-Z
    \end{enumerate}
    Equivalent matrix:
    \begin{equation*}
        Z H
    \end{equation*}
\end{examplebox}

\highspace
\begin{flushleft}
    \textcolor{Green3}{\faIcon{grip-lines} \textbf{Parallel Gates}}
\end{flushleft}
Now let's complicate things a bit. Consider this circuit with two qubits:
\begin{center}
    \begin{quantikz}
        \lstick{\ket{q_1}} & \gate{H} & \\
        \lstick{\ket{q_2}} & \gate{X} &
    \end{quantikz}
\end{center}

\noindent
The gates happen \textbf{at the same time}. There is no ``first'' or ``second'' gate. Instead, we combine them using the \textbf{tensor product} (\autopageref{subsubsec:tensor-product}). Thus, the equivalent matrix is:
\begin{equation*}
    H \otimes X
\end{equation*}
This is the operator acting on the two-qubit system.

\highspace
\begin{examplebox}
    Consider the following state: $\ket{\psi} = \ket{00}$. Apply two gates in parallel: Hadamard on the first qubit and Pauli-X on the second qubit. The circuit is:
    \begin{center}
        \begin{quantikz}
            \lstick{\ket{q_1}} & \gate{H} & \\
            \lstick{\ket{q_2}} & \gate{X} &
        \end{quantikz}
        =
        \begin{quantikz}
            \lstick{\ket{0}} & \gate{H} & \\
            \lstick{\ket{0}} & \gate{X} &
        \end{quantikz}
    \end{center}
    The equivalent matrix is $H \otimes X$. Applying this to the state $\ket{00}$ gives:
    \begin{align*}
        (H \otimes X) \ket{00} &= H\ket{0} \otimes X\ket{0} \\
        &= \frac{1}{\sqrt{2}} \left(\ket{0} + \ket{1}\right) \otimes \ket{1} \\
        &= \frac{1}{\sqrt{2}} \left(\ket{01} + \ket{11}\right)
    \end{align*}
\end{examplebox}

\highspace
In general, if we have a circuit with $n$ qubits and we apply $n$ gates in parallel, the equivalent matrix is the tensor product of all the gates:
\begin{equation*}
    U = U_1 \otimes U_2 \otimes \cdots \otimes U_n
\end{equation*}
Now suppose we have a circuit with \textbf{two layers}:
\begin{center}
    \begin{quantikz}[slice all]
        \lstick{\ket{q_1}} & \gate{H} & \gate{Z} & \\
        \lstick{\ket{q_2}} & \gate{X} & \gate{Y} &
    \end{quantikz}
\end{center}

\noindent
The first layer becomes $H \otimes X$, and the second layer becomes $Z \otimes Y$. The equivalent matrix for the entire circuit is:
\begin{equation*}
    (Z \otimes Y) \cdot (H \otimes X)
\end{equation*}
Again, notice that we still multiply \textbf{from right to left}, because the first layer acts first, and the second layer acts second.

\highspace
So there are \textbf{two different operations}:
\begin{itemize}
    \item \textbf{Horizontally (along time)}: matrix multiplication.
    \item \textbf{Vertically (same time, different qubits)}: tensor product.
\end{itemize}

\newpage

\begin{flushleft}
    \textcolor{Green3}{\faIcon{question-circle} \textbf{What if, vertically, one qubit has a gate and the other has a straight line?}}
\end{flushleft}
In this case, the scenario is really simple. Let's say we have a circuit with two qubits $\ket{\psi} = \ket{00}$, composed of:
\begin{center}
    \begin{quantikz}[slice all]
        \lstick{\ket{0}}\slice{0} & \gate{H} & \gate{Z} & \gate{S} & \\
        \lstick{\ket{0}}          & \gate{X} & \gate{Y} &          &
    \end{quantikz}
\end{center}

\noindent
We label each layer with a temporary state $\ket{\psi_i}$, where $i$ is the layer number:
\begin{itemize}
    \item \textbf{Initial state} $\ket{\psi_0}$. Both qubit starts in \ket{0}, therefore:
    \begin{equation*}
        \ket{\psi_0} = \ket{0} \otimes \ket{0} = \ket{00}
    \end{equation*}
    We apply the tensor product because the two qubits are vertically aligned.


    \item \textbf{First layer} $\ket{\psi_1}$. Since the gates are simultaneous, we apply the tensor product of the two gates:
    \begin{equation*}
        \ket{\psi_1} = (H \otimes X) \ket{\psi_0} = (H \otimes X) \ket{00}
    \end{equation*}
    Remember what we learned on \autopageref{ex:unitaries-pauli-gates-and-eigenstates-of-z-x-y-h}, instead of multiplying two huge matrices, we can apply the two gates separately to each qubit:
    \begin{align*}
        \ket{\psi_1} &= H\ket{0} \otimes X\ket{0} \\
        &= \underbrace{\frac{1}{\sqrt{2}} \left(\ket{0} + \ket{1}\right)}_{\ket{+}} \otimes \ket{1}
    \end{align*}
    \textcolor{Green3}{\faIcon[regular]{lightbulb} \textbf{Tip.}} At this point \textbf{do not expand the tensor product!} Leave it factorized, because it is easier to work with.


    \item \textbf{Second layer} $\ket{\psi_2}$. Again, we apply the tensor product of the two gates:
    \begin{equation*}
        \ket{\psi_2} = (Z \otimes Y) \ket{\psi_1} = (Z \otimes Y) \left[\frac{1}{\sqrt{2}} \left(\ket{0} + \ket{1}\right) \otimes \ket{1}\right]
    \end{equation*}
    \begin{itemize}
        \item First qubit (Pauli-Z changes only the relative phase of the \ket{1} qubit):
        \begin{equation*}
            Z \left[\frac{1}{\sqrt{2}} \left(\ket{0} + \ket{1}\right)\right] = \frac{1}{\sqrt{2}} \left(Z\ket{0} + Z\ket{1}\right) = \frac{1}{\sqrt{2}} \left(\ket{0} - \ket{1}\right)
        \end{equation*}
        \item Second qubit (Pauli-Y changes the relative phase of the \ket{1} qubit and flips it):
        \begin{equation*}
            Y \ket{1} = \begin{bmatrix}
                0 & -i \\
                i & 0
            \end{bmatrix} \begin{bmatrix}
                0 \\
                1
            \end{bmatrix} = \begin{bmatrix}
                0 + (-i) \cdot 1 \\
                0 + 0
            \end{bmatrix} =
            \begin{bmatrix}
                -i \\
                0
            \end{bmatrix} = -i \ket{0}
        \end{equation*}
    \end{itemize}
    Putting everything together:
    \begin{equation*}
        \ket{\psi_2} = \left[\frac{1}{\sqrt{2}} \left(\ket{0} - \ket{1}\right)\right] \otimes \left[-i \ket{0}\right] = -i \cdot \frac{1}{\sqrt{2}} \left(\ket{0} - \ket{1}\right) \otimes \ket{0}
    \end{equation*}
    \textcolor{Green3}{\faIcon{search} \textbf{Global Phase detected.}} Before we apply the last gate, we can see that the state has a term $-i$ in front of the tensor product that multiplies \textbf{the whole quantum state}. It is not attached only to one basis state (like $\ket{0} - \ket{1}$ on the first qubit). Therefore it has no physical effect and can be discarded. We can rewrite the state as:
    \begin{equation*}
        \ket{\psi_2} = \frac{1}{\sqrt{2}} \left(\ket{0} - \ket{1}\right) \otimes \ket{0}
    \end{equation*}


    \item \textbf{Third layer} $\ket{\psi_3}$. The last layer has a gate on the first qubit and nothing on the second qubit. In this case, we apply the \textbf{tensor product of the gate} and the \textbf{identity matrix}:
    \begin{equation*}
        \ket{\psi_3} = (S \otimes I) \ket{\psi_2} = (S \otimes I) \left[\frac{1}{\sqrt{2}} \left(\ket{0} - \ket{1}\right) \otimes \ket{0}\right]
    \end{equation*}
    \begin{itemize}
        \item First qubit (phase gate $S$ changes only the relative phase of the \ket{1} qubit):
        \begin{equation*}
            S \left[\frac{1}{\sqrt{2}} \left(\ket{0} - \ket{1}\right)\right] = \frac{1}{\sqrt{2}} \left(S\ket{0} - S\ket{1}\right) = \frac{1}{\sqrt{2}} \left(\ket{0} - i \ket{1}\right)
        \end{equation*}
        \item Second qubit (Identity gate does not change the state):
        \begin{equation*}
            I \ket{0} = \begin{bmatrix}
                1 & 0 \\
                0 & 1
            \end{bmatrix} \begin{bmatrix}
                1 \\
                0
            \end{bmatrix} = \begin{bmatrix}
                1 \\
                0
            \end{bmatrix} = \ket{0}
        \end{equation*}
    \end{itemize}
    Remembering that we can apply the two gates separately to each qubit, we can write the final state as:
    \begin{equation*}
        \ket{\psi_3} = \left[\frac{1}{\sqrt{2}} \left(\ket{0} - i \ket{1}\right)\right] \otimes \ket{0}
    \end{equation*}
\end{itemize}
Now we are at the end of the circuit, so we can \textbf{expand the tensor product} to get the final state:
\begin{equation*}
    \ket{\psi} = \frac{1}{\sqrt{2}} \left(\ket{00} - i \ket{10}\right)
\end{equation*}
From this final state, we can calculate the measurement probabilities and the marginal probabilities of each qubit:
\begin{itemize}
    \item Measurement probabilities:
    \begin{align*}
        P(00) &= \left|\frac{1}{\sqrt{2}}\right|^2 = \frac{1}{2} \\
        P(01) &= 0 \\
        P(10) &= \left|-\frac{i}{\sqrt{2}}\right|^2 = \left|-\frac{1}{\sqrt{2}} \cdot i\right|^{2} = \left(\sqrt{0^2 + \left(-\frac{1}{\sqrt{2}}\right)^2}\right)^{2} = \frac{1}{2} \\
        P(11) &= 0
    \end{align*}
    \item Marginal probabilities:
    \begin{align*}
        P(q_1 = 0) &= P(00) + P(01) = \frac{1}{2} + 0 = \frac{1}{2} \\
        P(q_1 = 1) &= P(10) + P(11) = \frac{1}{2} + 0 = \frac{1}{2} \\
        P(q_2 = 0) &= P(00) + P(10) = \frac{1}{2} + \frac{1}{2} = 1 \\
        P(q_2 = 1) &= P(01) + P(11) = 0 + 0 = 0
    \end{align*}
\end{itemize}

\begin{flushleft}
    \textcolor{Green3}{\faIcon{question-circle} \textbf{What if we have a circuit with more than two qubits?}}
\end{flushleft}
The workflow does not change when the number of qubits increases. The only difference is that the tensor product will have more terms, and the equivalent matrix will be larger. The same rules apply: multiply matrices from right to left for consecutive gates, and use the tensor product for parallel gates. Consider the following circuit with three qubits:
\begin{center}
    \begin{quantikz}[slice all]
        \lstick{\ket{0}}\slice{0} & \gate{H} & \gate{Y} & \\
        \lstick{\ket{1}}          & \gate{H} & \gate{Z} & \\
        \lstick{\ket{0}}          & \gate{X} & \gate{S} & 
    \end{quantikz}
\end{center}

\noindent
There are three layers, and the equivalent matrix is:
\begin{equation*}
    (Y \otimes Z \otimes S) \cdot (H \otimes H \otimes X) \cdot \left(\ket{0} \otimes \ket{1} \otimes \ket{0}\right)
\end{equation*}
Let's label each layer with a temporary state $\ket{\psi_i}$, where $i$ is the layer number:
\begin{itemize}
    \item \textbf{Initial state} $\ket{\psi_0}$. The three qubits start in $\ket{010}$, therefore:
    \begin{equation*}
        \ket{\psi_0} = \ket{0} \otimes \ket{1} \otimes \ket{0} = \ket{010}
    \end{equation*}
    The order is always from top to bottom, because the first qubit is the most significant bit and the last qubit is the least significant bit: $\ket{q_1 q_2 q_3}$.


    \item \textbf{First layer} $\ket{\psi_1}$. Since the gates are simultaneous, we apply the tensor product of the three gates:
    \begin{equation*}
        \begin{aligned}
            \ket{\psi_1} &= (H \otimes H \otimes X) \ket{\psi_0} \\
            &= (H \otimes H \otimes X) \ket{010} \\
            &= H\ket{0} \otimes H\ket{1} \otimes X\ket{0} \\
            &= \frac{1}{\sqrt{2}} \left(\ket{0} + \ket{1}\right) \otimes \frac{1}{\sqrt{2}} \left(\ket{0} - \ket{1}\right) \otimes \ket{1}
        \end{aligned}
    \end{equation*}
    \textcolor{Green3}{\faIcon[regular]{lightbulb}} Again, leave the tensor product factorized, because it is easier to work with.


    \item \textbf{Second layer} $\ket{\psi_2}$. Again, we apply the tensor product of the three gates:
    \begin{equation*}
        \begin{aligned}
            \ket{\psi_2} &= (Y \otimes Z \otimes S) \ket{\psi_1} \\
            &= (Y \otimes Z \otimes S) \left[\frac{1}{\sqrt{2}} \left(\ket{0} + \ket{1}\right) \otimes \frac{1}{\sqrt{2}} \left(\ket{0} - \ket{1}\right) \otimes \ket{1}\right]
        \end{aligned}
    \end{equation*}
    \begin{itemize}
        \item First qubit (Pauli-Y changes the relative phase of the \ket{1} qubit and flips it):
        \begin{equation*}
            Y \left[\frac{1}{\sqrt{2}} \left(\ket{0} + \ket{1}\right)\right] = \frac{1}{\sqrt{2}} \left(Y\ket{0} + Y\ket{1}\right)
        \end{equation*}
        \begin{align*}
            Y\ket{0} &= \begin{bmatrix}
                0 & -i \\
                i & 0
            \end{bmatrix} \begin{bmatrix}
                1 \\
                0
            \end{bmatrix} = \begin{bmatrix}
                0 \\
                i
            \end{bmatrix} = i \ket{1} \\
            Y\ket{1} &= \begin{bmatrix}
                0 & -i \\
                i & 0
            \end{bmatrix} \begin{bmatrix}
                0 \\
                1
            \end{bmatrix} = \begin{bmatrix}
                -i \\
                0
            \end{bmatrix} = -i \ket{0}
        \end{align*}
        Therefore:
        \begin{equation*}
            Y \left[\frac{1}{\sqrt{2}} \left(\ket{0} + \ket{1}\right)\right] = \frac{1}{\sqrt{2}} \left(-i \ket{0} + i \ket{1}\right) = i \cdot \frac{1}{\sqrt{2}} \left(\ket{1} - \ket{0}\right)
        \end{equation*}
        We don't like how the minus sign is in front of the \ket{0} term, so we can factor out a $-1$ to get:
        \begin{equation*}
            i \cdot \frac{1}{\sqrt{2}} \left(\ket{1} - \ket{0}\right) = -i \cdot \frac{1}{\sqrt{2}} \left(\ket{0} - \ket{1}\right)
        \end{equation*}
        \item Second qubit (Pauli-Z changes only the relative phase of the \ket{1} qubit):
        \begin{equation*}
            Z \left[\frac{1}{\sqrt{2}} \left(\ket{0} - \ket{1}\right)\right] = \frac{1}{\sqrt{2}} \left(Z\ket{0} - Z\ket{1}\right) = \frac{1}{\sqrt{2}} \left(\ket{0} + \ket{1}\right)
        \end{equation*}
        \item Third qubit (phase gate $S$ changes only the relative phase of the \ket{1} qubit):
        \begin{equation*}
            S \ket{1} = i \ket{1}
        \end{equation*}
    \end{itemize}
    The final state is:
    \begin{equation*}
        \ket{\psi_2} = \left[-i \cdot \frac{1}{\sqrt{2}} \left(\ket{0} - \ket{1}\right)\right] \otimes \left[\frac{1}{\sqrt{2}} \left(\ket{0} + \ket{1}\right)\right] \otimes \left(i \ket{1}\right)
    \end{equation*}
\end{itemize}
We can simplify the global phase $-i$ and $i$: $-i \cdot i = -i^2 = 1$. Therefore, the final state is:
\begin{equation*}
    \ket{\psi} = \frac{1}{\sqrt{2}} \left(\ket{0} - \ket{1}\right) \otimes \frac{1}{\sqrt{2}} \left(\ket{0} + \ket{1}\right) \otimes \ket{1}
\end{equation*}
We can also simplify the tensor product of the first two qubits and leave the third qubit factorized:
\begin{align*}
    \ket{\psi} &= \frac{1}{2} \left(\ket{0}\ket{0} + \ket{0}\ket{1} - \ket{1}\ket{0} - \ket{1}\ket{1}\right) \otimes \ket{1} \\
    &= \frac{1}{2} \left(\ket{00} + \ket{01} - \ket{10} - \ket{11}\right) \otimes \ket{1} \\
    &= \frac{1}{2} \left(\ket{001} + \ket{011} - \ket{101} - \ket{111}\right)
\end{align*}